\section{Discuss\~ao}
\label{sec:discussao}

\subsection{O que os dois m\'etodos oferecem, pelo que foi medido}

\begin{center}
\begin{adjustbox}{max width=\linewidth}
\begin{tabular}{p{4.6cm}p{4.6cm}p{4.6cm}}
\toprule
 & front-tracking (balanceado) & VOF \\
\midrule
curvatura / salto de press\~ao & exata no c\'irculo; erro \num{1.7e-7}
  & derivada de campo suavizado; erro \num{4.7e-3} \\
correntes par\'asitas & $<\num{5e-11}$, decaindo & \num{5.2e-5}, decaindo \\
massa & conservada enquanto o escoamento \'e est\'atico; deriva com o movimento
  & conservada por constru\c{c}\~ao em malha FIXA; com remalhamento adaptativo
    deriva \num{1e-3}, como o front-tracking (Se\c{c}\~ao~\ref{sec:massa-amr}) \\
topologia (coalesc\^encia, quebra) & manual, fora do escopo & autom\'atica \\
custo de estrutura & marcadores ordenados, cirurgia, indicadora por c\'elula
  & campo euleriano a mais \\
\bottomrule
\end{tabular}
\end{adjustbox}
\end{center}

A leitura n\~ao \'e ``um vence''. Nas m\'etricas deste teste o front-tracking
balanceado sai \`a frente, mas o teste \'e uma gota est\'atica: n\~ao exercita
topologia, n\~ao exercita deforma\c{c}\~ao grande, e n\~ao exercita a
conserva\c{c}\~ao de massa sob movimento --- que \'e justamente onde o VOF tem
vantagem \emph{estrutural}, e n\~ao circunstancial. Escoamento em que a massa da
fase precisa ser conservada ao longo de muitos tempos, ou em que gotas coalescem
e se partem, continua favorecendo o VOF. Com os dois no mesmo sistema, a escolha
passa a ser informada por medida em vez de por prefer\^encia.

\subsection{A compara\c{c}\~ao como verifica\c{c}\~ao m\'utua}

Vale sublinhar um uso da compara\c{c}\~ao que n\~ao \'e o de escolher m\'etodo.
Como a dire\c{c}\~ao de v\'arios resultados \'e conhecida \emph{a priori} pela
literatura, eles funcionam como or\'aculo um para o outro: se o VOF tivesse
sa\'ido melhor em curvatura, ou se o front-tracking espalhado tivesse conservado
massa exatamente, o resultado apontaria defeito de implementa\c{c}\~ao, n\~ao
descoberta sobre os m\'etodos.

E foi exatamente assim que o defeito apareceu. O \'unico resultado que
\emph{contrariou} a expectativa --- correntes par\'asitas maiores no
front-tracking --- n\~ao era propriedade do m\'etodo: era uma deficiência da
implementa\c{c}\~ao, localizada pela discrep\^ancia e removida na
Se\c{c}\~ao~\ref{sec:balanco}. Sem o VOF ao lado, com o mesmo problema e as
mesmas sondas, n\~ao haveria com o que estranhar \num{5e-4}: pareceria apenas
``pequeno''.

\subsection{Limita\c{c}\~oes}

Para que o alcance dos n\'umeros n\~ao seja superestimado:

\begin{itemize}
  \item \textbf{Uma \'unica resolu\c{c}\~ao.} Tudo foi medido a $10{,}2$
    c\'elulas por raio. Um estudo de converg\^encia --- refinar $h$ e obter a
    taxa de cada m\'etrica --- n\~ao foi feito.
  \item \textbf{A vantagem do balan\c{c}o \'e espec\'ifica do caso est\'atico
    com $\kappa$ constante}, como a Se\c{c}\~ao~\ref{sec:elipse} mostrou: na
    elipse as duas formula\c{c}\~oes ficam equivalentes dentro da medida.
  \item \textbf{A curvatura na faceta \'e uma m\'edia ponderada dos marcadores
    vizinhos}, e n\~ao um campo de curvatura propriamente constru\'ido e
    estendido. Foi o suficiente para estabilizar a elipse, mas n\~ao foi testada
    em geometria com curvatura fortemente vari\'avel ou com mudan\c{c}a de
    sinal.
  \item \textbf{Serial nas medidas, mas n\~ao mais por impedimento.} Todos os
    n\'umeros deste relat\'orio s\~ao $n_p=1$. O que mudou \'e que $n_p=2$ foi
    testado e \emph{consertado} onde precisava: as subse\c{c}\~oes
    \ref{sec:paralelo} e \ref{sec:posse} tra\c{c}am o que quebrava, o que foi
    corrigido e o que se revelou nunca ter estado quebrado. Falta rodar em
    paralelo um caso longo e refazer as medidas --- o que nao foi feito.
  \item \textbf{Escalonamento: n\~ao medido, e a arquitetura atual o limita.}
    Nenhum estudo de escalabilidade foi feito --- n\~ao havia como, porque at\'e
    este trabalho o remalhamento estava errado em $n_p>1$ e o caso o exige. Mas
    tr\^es fatos do c\'odigo, esses \emph{medidos}, dizem onde o limite estaria:
    \begin{enumerate}
      \item o caso tem \textbf{uma \'unica \'arvore}
        (\texttt{number\_domains:\ 1}), e o la\c{c}o de leitura reparte
        \emph{arquivos} entre \emph{ranks}: com um arquivo s\'o, o rank 0
        materializa a malha global inteira em mem\'oria antes de o balanceador a
        repartir. \'E um teto de mem\'oria por rank que n\~ao melhora com mais
        processos;
      \item \texttt{higflow\_\allowbreak{}reconstroi\_\allowbreak{}dominio} \textbf{reconstr\'oi o dom\'inio
        do zero} a cada remalhamento --- \texttt{create\_domain} mais
        \texttt{initialize\_\allowbreak{}domain\_\allowbreak{}yaml} --- em vez de adaptar no lugar;
      \item o crit\'erio passa por \textbf{arquivo, escrito por um rank s\'o},
        com barreira, e relido por todos (Se\c{c}\~ao~\ref{sec:paralelo}). Foi o
        conserto correto para a \emph{corre\c{c}\~ao}, e \'e um trecho serial
        proporcional \`a malha global a cada remalhamento.
    \end{enumerate}
    A consequ\^encia pr\'atica, e ela contraria a leitura natural: usar o t8code
    apenas como \emph{particionador} --- que \'e o caminho barato, porque
    \texttt{fonte\_de\_particao} j\'a existe --- melhoraria a qualidade da
    parti\c{c}\~ao e \textbf{n\~ao} entregaria o escalonamento, que \'e a raz\~ao
    de ser do t8code. Os tr\^es itens acima continuariam intactos. O que
    entregaria \'e a adapta\c{c}\~ao din\^amica distribu\'ida: crit\'erio avaliado
    localmente, floresta adaptada no lugar, e o \texttt{.amr} reduzido a formato
    de \emph{arranque} em vez de via de remalhamento. Isso \'e reescrever o
    caminho de remalhamento, n\~ao ligar um gancho.

    \textbf{Isso foi medido, e a infer\^encia caiu.} A
    Se\c{c}\~ao~\ref{sec:fase0} cronometra o remalhamento separado do resto em 1,
    2, 4 e 8 processos: ele custa entre \SI{0.06}{\percent} e \SI{0.10}{\percent}
    do tempo, e o trecho serial da escrita do \texttt{.amr} \'e \SI{3}{\percent}
    \emph{disso}. Os tr\^es gargalos s\~ao reais como fatos de c\'odigo e
    irrelevantes como custo, neste tamanho de problema. O que a medida encontrou
    no lugar \'e que a efici\^encia paralela cai a \SI{22}{\percent} em oito
    processos --- o caso \'e pequeno demais para o paralelismo render.
  \item \textbf{t8code: integrado no monof\'asico, n\~ao no bif\'asico.} O gancho
    \texttt{malha\_t8\_instala} existe nos dois exemplos e instala a fonte, mas
    \texttt{higflow\_\allowbreak{}partition\_\allowbreak{}domain\_\allowbreak{}multiphase} n\~ao a consulta --- l\^e o
    \texttt{.amr} sempre. A fronteira imersa funciona com t8code porque \'e
    monof\'asica e passa por outro caminho. Desde este trabalho o pedido
    ignorado \'e ao menos \emph{denunciado} na linha de proced\^encia da malha, em
    vez de passar em sil\^encio. E h\'a um segundo impedimento atr\'as do
    primeiro: as fontes do t8code recusam \texttt{.amr} multin\'ivel
    (\texttt{numlevels != 1}, com a mensagem ``recusar \'e melhor que
    aproximar''), e o crit\'erio adaptativo escreve exatamente isso. Integrar de
    verdade exige uma fonte que \emph{leia} a lista de refino --- que o
    \texttt{higio\_amr\_info} carrega e entrega, e que hoje \'e descartada em
    favor s\'o da caixa.
  \item \textbf{Propriedades iguais.} A terceira mudan\c{c}a estrutural --- o
    salto de $\rho$ e $\mu$ --- est\'a adiada nos dois lados. Em escoamento com
    raz\~ao de densidade alta as conclus\~oes podem mudar.
  \item \textbf{A cirurgia foi exercitada} apenas na relaxa\c{c}\~ao da elipse,
    onde a frente encurta monotonicamente (108 para 102 marcadores). Esticamento
    forte, que exigiria \emph{inser\c{c}\~ao} sustentada, segue sem teste
    acoplado.
\end{itemize}

\subsection{Paralelo: o que foi medido em $n_p=2$}
\label{sec:paralelo}

Este relat\'orio dizia, at\'e aqui, apenas que o trabalho era serial. Isso era
verdade e insuficiente: ``n\~ao implementado'' e ``implementado e errado'' pedem
provid\^encias diferentes, e a diferen\c{c}a s\'o se conhece rodando. Rodou-se,
em \'arvore de trabalho isolada, o mesmo caso curto com $n_p=1$ e $n_p=2$ nos
dois m\'etodos.

\begin{adjustbox}{max width=\linewidth}
\pgfplotstabletypeset[
  string type,
  columns/grandeza/.style={column name={grandeza}},
  columns/np1/.style={column name={$n_p=1$}},
  columns/np2/.style={column name={$n_p=2$}},
  every head row/.style={before row=\toprule, after row=\midrule},
  every last row/.style={after row=\bottomrule},
]{\dat{paralelo.dat}}
\end{adjustbox}

\paragraph{O front-tracking congela quase metade da frente.} Em $n_p=2$,
\textbf{109 dos 252 marcadores n\~ao se moveram} em 10 passos, e formam uma faixa
\emph{cont\'igua} em $y$ --- a parte da bolha que cai no subdom\'inio do outro
processo. A causa \'e uma linha da interpola\c{c}\~ao de velocidade, que em s\'erie
nunca dispara:
\begin{center}
\texttt{if (h <= 0.0) return;\qquad // fora do dominio: marcador nao anda}
\end{center}
A frente \'e \emph{replicada}: cada processo advecta a sua pr\'opria c\'opia com
a velocidade que enxerga, e as c\'opias divergem no primeiro passo. A prova
direta \'e que os dois processos relatam desvios de parti\c{c}\~ao da unidade
\emph{diferentes} ($0{,}402$ e $0{,}795$) --- n\~ao est\~ao vendo o mesmo
est\^encil. Contra \num{3.05e-14} em s\'erie, \'e a mesma assinatura que
identificou a ``for\c{c}a exatamente pela metade'' no corpo r\'igido: pesos que
n\~ao somam 1 devolvem velocidade pequena demais, sem erro vis\'ivel.

\paragraph{O VOF acerta o solver e erra o remalhamento.} Sem remalhar, as
s\'eries de $n_p=1$ e $n_p=2$ saem \textbf{id\^enticas em todos os d\'igitos
impressos}: o caminho multif\'asico do solver e as m\'etricas (que reduzem com
\texttt{MPI\_Allreduce}) est\~ao corretos em paralelo. Com remalhamento, n\~ao:
o pr\'oprio contador do remalhamento acusa \textbf{265 posi\c{c}\~oes sem valor}
no primeiro remalhamento e \textbf{2709} no segundo, e $V_c$ sai com
\SI{77}{\percent} de erro --- \emph{caindo} quando devia subir.

Vale notar o que esse contador \'e: o n\'umero de posi\c{c}\~oes da malha nova
para as quais a colheita n\~ao tinha valor. A mensagem que o imprime existia nos
dois exemplos e vinha imprimindo zero desde sempre, porque em s\'erie \'e zero.
Era um or\'aculo instalado que nunca havia disparado.

\paragraph{Os consertos, e o que aconteceu com cada um.}\mbox{}

\begin{adjustbox}{max width=\linewidth}
\pgfplotstabletypeset[
  string type,
  columns/conserto/.style={column name={conserto}},
  columns/estado/.style={column name={estado}},
  columns/evidencia/.style={column name={evid\^encia}},
  every head row/.style={before row=\toprule, after row=\midrule},
  every last row/.style={after row=\bottomrule},
]{\dat{paralelo-consertos.dat}}
\end{adjustbox}

Dois s\~ao mec\^anicos: os escritores de crit\'erio faziam
\texttt{fopen(...,"w")} do \emph{mesmo} caminho em todos os processos, sem guarda
de \emph{rank} e sem barreira antes da releitura --- e os dois s\~ao necess\'arios,
porque a guarda sozinha n\~ao faz ningu\'em esperar.

O terceiro n\~ao se cura com guarda: as sementes do VOF saem do dom\'inio
\emph{local} sem \emph{gather}, de modo que mesmo o processo 0 escreveria uma
malha refinada s\'o em torno do peda\c{c}o de interface que ele possui. A
evid\^encia de que o \emph{gather} resolveu veio da pr\'opria contagem de
sementes: antes, os dois processos relatavam $122+186=308$ num remalhamento e
$150+158=308$ no seguinte; depois, ambos relatam 308. A soma das partes antes \'e
exatamente o todo depois. (No front-tracking esse defeito n\~ao existe: a frente
\'e replicada, ent\~ao todo processo tem a geometria completa.)

O quarto era o \'unico sem c\'odigo pronto para copiar, e est\'a descrito na
subse\c{c}\~ao seguinte.

\subsection{A advec\c{c}\~ao paralela, e uma analogia que enganou nos dois sentidos}
\label{sec:posse}

A velocidade de um marcador \emph{n\~ao \'e calcul\'avel localmente}: o suporte do
n\'ucleo pode estar repartido entre processos. Tr\^es desenhos foram considerados,
e a parti\c{c}\~ao da unidade decidiu entre eles por medida:

\begin{itemize}
  \item \textbf{somar os parciais} --- duplica. \texttt{\_suporte} localiza
    facetas no dom\'inio local, que inclui \emph{franja}: 172 das 504 entradas
    eram encontradas \textbf{completas} pelos dois processos, e a soma daria o
    dobro exato.
  \item \textbf{normalizar pela soma dos pesos} --- n\~ao serve. 12 entradas
    tinham um processo \textbf{completo} e outro \textbf{parcial}; dividir mistura
    a interpola\c{c}\~ao boa com uma estimativa parcial diferente.
  \item \textbf{filtro de posse} --- contribui apenas o processo cuja soma
    \emph{local} de pesos \'e 1, e divide pela \emph{contagem} de processos
    completos. Exato, porque eles somam sobre o \emph{mesmo} conjunto de facetas.
\end{itemize}

O resultado: os marcadores congelados passam de 109 a \textbf{zero}, e a frente de
$n_p=2$ fica a \num{1.5e-11} da serial --- abaixo da toler\^ancia do solver linear
(\num{1e-10}), contra \num{2.3e-6} do defeito. Em $n_p=1$ o caminho permanece
\textbf{id\^entico byte a byte}, que \'e a condi\c{c}\~ao para as corridas j\'a
feitas seguirem compar\'aveis.

A medida tamb\'em mostrou que o filtro \'e sempre aplic\'avel aqui --- nenhuma
entrada ficou sem ao menos um processo completo. Onde isso falhar (mais processos,
outra parti\c{c}\~ao) n\~ao existe resposta certa a dar, e o c\'odigo
\textbf{aborta}: devolver zero seria repor o defeito original com outra roupa.

\paragraph{O quinto conserto n\~ao existia.} Este relat\'orio afirmou, numa vers\~ao
anterior desta se\c{c}\~ao, que o \emph{espalhamento} do front-tracking precisava
da redu\c{c}\~ao franja$\to$dono que \texttt{fi\_espalha} faz vinte linhas adiante.
Medido, cai:

\begin{adjustbox}{max width=\linewidth}
\pgfplotstabletypeset[
  string type,
  columns/regra-de-posse-no-espalhamento/.style={column name={regra de posse no espalhamento}},
  columns/np1/.style={column name={$n_p=1$}},
  columns/np2contraserial/.style={column name={$n_p=2$ vs serial}},
  every head row/.style={before row=\toprule, after row=\midrule},
  every last row/.style={after row=\bottomrule},
]{\dat{paralelo-posse.dat}}
\end{adjustbox}

A regra original j\'a estava correta, e a redu\c{c}\~ao era a \'unica das tr\^es
que \emph{quebrava}. O motivo \'e que o dono de uma faceta acumula de todos os
marcadores que consegue localizar, e todo marcador cujo suporte toca aquela faceta
est\'a a no m\'aximo $1{,}5$ c\'elula dela --- portanto dentro do dom\'inio dele
\emph{com franja}. A acumula\c{c}\~ao do dono j\'a \'e completa, e o
\texttt{dp\_sync} descartar o que os outros escreveram na franja \'e precisamente
o que torna o espalhamento de frente \textbf{replicada} correto sem
comunica\c{c}\~ao nenhuma.

A analogia com o corpo r\'igido enganou nos \emph{dois} sentidos: fez ver defeito
onde n\~ao havia --- porque l\'a a redu\c{c}\~ao \'e necess\'aria --- e ofereceu a
cura errada, porque l\'a o \texttt{DMSwarm} reparte os marcadores por posse
exclusiva e aqui a frente \'e replicada. Mesmo arquivo, modelos de
distribui\c{c}\~ao incompat\'iveis.

\paragraph{E a conserva\c{c}\~ao n\~ao decidiu nada.} As tr\^es regras produziram
n\'umeros de conserva\c{c}\~ao \emph{id\^enticos em todos os d\'igitos} ---
inclusive a que degrada a frente em cinco ordens de grandeza. Numa gota em repouso
os dois lados s\~ao \num{4.4e-12}: zero contra zero. A soma n\~ao v\^e \emph{onde}
a for\c{c}a foi posta, e decidir por ela teria escolhido qualquer uma das tr\^es.
\'E a mesma li\c{c}\~ao da Se\c{c}\~ao~\ref{sec:verificacao}, onde a
conserva\c{c}\~ao passava e a parti\c{c}\~ao da unidade acusava metade da
for\c{c}a.

\paragraph{Uma armadilha de medi\c{c}\~ao, de gra\c{c}a.} O caso com remalhamento
rodou em \SI{64}{\second} em $n_p=2$ contra \SI{170}{\second} em $n_p=1$: $2{,}7$
vezes em dois processos, \emph{superlinear}. A superlinearidade \'e a assinatura
de resolver um problema \emph{menor}, e n\~ao mais r\'apido. Quem medisse
escalabilidade neste estado veria um resultado excelente.

\paragraph{E o que os tr\^es testes t\^em em comum.} Nenhum quebrou. Tr\^es
execu\c{c}\~oes com c\'odigo de sa\'ida zero, resultados de aspecto plaus\'ivel,
e defeitos de \SI{43}{\percent}, \SI{77}{\percent} e \SI{100}{\percent}. O que os
separou do correto foram os or\'aculos --- um deles j\'a instalado e nunca
disparado.

\subsection{Pr\'oximos passos}

Em ordem de valor:

\begin{enumerate}
  \item \textbf{Levar a s\'erie a uma malha mais fina que $40$ c\'elulas por
    raio.} A s\'erie com banda constante ficou mon\'otona nas cinco m\'etricas
    (Se\c{c}\~ao~\ref{sec:conv-serieC}), mas em tr\^es delas a diferen\c{c}a
    entre malhas sucessivas ainda \emph{cresce} ao refinar: n\~ao h\'a regime
    assint\'otico, e ordem continua sem poder ser declarada. \'E o pr\'oximo
    ponto que decide, e custa $4$ a $8$ vezes a corrida de $40$.
  \item \textbf{Fase B3}, oscila\c{c}\~ao de gota contra a frequ\^encia de Lamb:
    exercita din\^amica acoplada e, por deformar a interface, finalmente p\~oe a
    cirurgia \`a prova com o solver no circuito.
  \item \textbf{Paralelo: repetir as medidas.} Os consertos das
    subse\c{c}\~oes~\ref{sec:paralelo} e \ref{sec:posse} foram verificados em
    casos de dez passos, com $n_p=2$ e uma \'unica parti\c{c}\~ao. Falta refazer
    as compara\c{c}\~oes deste relat\'orio em corrida longa e com mais
    processos --- e em particular conferir se o filtro de posse continua
    encontrando processo completo para todo marcador, que \'e a condi\c{c}\~ao
    sob a qual ele aborta.
  \item \textbf{Converg\^encia em $h$, e agora por raz\~ao medida.} O Hysing
    fechou (Se\c{c}\~ao~\ref{sec:fechamento}) e os cinco desvios s\~ao
    \emph{negativos}: h\'a vi\'es compartilhado pelos dois m\'etodos, e os tr\^es
    candidatos --- resolu\c{c}\~ao \'unica, restri\c{c}\~ao capilar de passo,
    refino da esteira --- s\'o se separam refinando. Isto deixa de ser item de
    completude e passa a ser a pergunta aberta do trabalho.

    A s\'erie com banda constante j\'a estreitou um dos tr\^es: o refino da
    esteira \emph{importa}, e importa por largura absoluta e n\~ao por
    const\^ancia --- a s\'erie cuja banda nunca desce de $0{,}0625$ fica
    mon\'otona, a que desce a $0{,}0312$ n\~ao. Mas os desvios continuam todos
    negativos na s\'erie nova, e a mais fina \'e a que mais se afasta em
    $V_{c,\max}$. O vi\'es compartilhado segue de p\'e.
\end{enumerate}

A compara\c{c}\~ao com o VOF pode acompanhar cada uma dessas etapas: o caso
existe, as sondas est\~ao instaladas nos dois lados, e as duas formula\c{c}\~oes
de for\c{c}a coexistem atr\'as de um seletor.
